\section{Vigas sísmicas}
 
\subsection{Geometría y ubicación}
 
Las vigas sísmicas B7, B8 y B39 son elementos de hormigón armado que forman parte del sistema
resistente a cargas laterales del edificio. Los esfuerzos de diseño se obtienen de la modelación en
SAP2000 a partir de la envolvente de combinaciones de carga última sobre los 20 pisos del edificio.
La geometría varía en función del nivel: en pisos 1 al 3 se emplea ancho $b = 0{,}25$ m, mientras
que en pisos 4 al 20 se adopta $b = 0{,}20$ m. La altura total es constante por tipo de viga en
toda la altura del edificio. Las dimensiones se resumen en la Tabla~\ref{tab:viga_geometria}.
 
\begin{table}[H]
\centering
\caption{Geometría de las vigas sísmicas por zona del edificio.}
\label{tab:viga_geometria}
\begin{tabular}{lcccc}
\hline
Viga & $h$ (m) & $b$ (m), pisos 1--3 & $b$ (m), pisos 4--20 & $d$ (m) \\
\hline
B7  & 0,74 & 0,25 & 0,20 & 0,69 \\
B8  & 0,74 & 0,25 & 0,20 & 0,69 \\
B39 & 1,64 & 0,25 & 0,20 & 1,59 \\
\hline
\end{tabular}
\end{table}
 
La altura útil se calcula como $d = h - r_\text{ec}$, con recubrimiento $r_\text{ec} = 0{,}05$ m.
 
\subsection{Materiales}
 
Las vigas sísmicas emplean los mismos materiales que el resto de los elementos del edificio en los
niveles correspondientes:
 
Hormigón G35 en pisos 1 al 3 ($f_c' = 35\;\text{MPa}$) y G25 en pisos 4 en adelante ($f_c' = 25\;\text{MPa}$), especificados según NCh170.
 
Acero de refuerzo A630-420H según NCh204, con $f_y = 420\;\text{MPa}$ y $E_s = 200\;\text{GPa}$.
 
Los factores de reducción de resistencia aplicados son $\varphi_\text{flex} = 0{,}90$ y
$\varphi_\text{corte} = 0{,}75$, conforme a ACI 318-19.
 
\subsection{Esfuerzos de diseño}
 
Los esfuerzos de diseño corresponden a la envolvente de valores máximos absolutos de momento
flector $M_3$ y cortante $V_2$ entre todos los casos de carga del modelo:
\begin{equation}
    M_u = \max\!\left(|M_{3,\text{máx}}^+|,\;|M_{3,\text{mín}}^-|\right)
\end{equation}
\begin{equation}
    V_u = \max\!\left(|V_{2,\text{máx}}^+|,\;|V_{2,\text{mín}}^-|\right)
\end{equation}
 
Los esfuerzos aumentan progresivamente desde los pisos superiores hacia los inferiores, con los
valores más críticos concentrados en los pisos 1 al 3. Los esfuerzos máximos por viga se presentan
en la Tabla~\ref{tab:viga_esfuerzos}.
 
\begin{table}[H]
\centering
\caption{Esfuerzos de diseño máximos por viga (piso crítico = piso 1).}
\label{tab:viga_esfuerzos}
\begin{tabular}{lccc}
\hline
Viga & Piso crítico & $M_u$ (tonf$\cdot$m) & $V_u$ (tonf) \\
\hline
B7  & 1 & 23,43 & 28,13 \\
B8  & 1 & 18,75 & 50,68 \\
B39 & 1 & 29,21 & 77,79 \\
\hline
\end{tabular}
\end{table}
 
La viga B39 en piso 1 es el elemento más solicitado de la estructura, con $V_u = 77{,}79$ tonf y
$M_u = 29{,}21$ tonf$\cdot$m.
 
\subsection{Diseño a flexión}
 
\subsubsection{Armadura mínima y máxima}
 
Las cuantías límite se determinan conforme a ACI 318-19 (sección 9.6.1.2):
\begin{equation}
    A_{s,\text{mín}} = \max\!\left(\frac{0{,}25\sqrt{f_c'}}{f_y}\,b\,d\;,\;\frac{1{,}4}{f_y}\,b\,d\right)
\end{equation}
 
La profundidad máxima del bloque de compresión corresponde al límite de zona controlada por
tracción ($\varepsilon_s \geq 0{,}004$):
\begin{equation}
    c_{\text{máx}} = \frac{\varepsilon_{cu}}{\varepsilon_{cu} + 0{,}004}\,d = 0{,}375\,d
    \qquad \Rightarrow \qquad
    a_{\text{máx}} = \beta_1\,c_{\text{máx}}
\end{equation}
\begin{equation}
    A_{s,\text{máx}} = \frac{0{,}85\,f_c'\,a_{\text{máx}}\,b}{f_y}
\end{equation}
 
con $\beta_1 = 0{,}80$ para $f_c' > 28\;\text{MPa}$ y $\beta_1 = 0{,}85$ para $f_c' \leq 28\;\text{MPa}$.
 
\subsubsection{Armadura adoptada y verificaciones}
 
El área de acero requerida se obtiene a partir del equilibrio de la sección fisurada:
\begin{equation}
    a = \frac{A_s\,f_y}{0{,}85\,f_c'\,b}
    \qquad
    c = \frac{a}{\beta_1}
    \qquad
    \varepsilon_s = \varepsilon_{cu}\,\frac{d - c}{c} \geq \varepsilon_y
\end{equation}
\begin{equation}
    \varphi M_n = \varphi\,A_s\,f_y\!\left(d - \frac{a}{2}\right) \geq M_u
\end{equation}
 
Los resultados del diseño a flexión se presentan en la Tabla~\ref{tab:viga_flexion}. La armadura
adoptada en cada zona cumple los límites de cuantía mínima y máxima, y la sección opera en
régimen controlado por tracción en todos los casos.
 
\begin{table}[H]
\centering
\caption{Armadura longitudinal adoptada por zona del edificio.}
\label{tab:viga_flexion}
\begin{tabular}{llccc}
\hline
Zona (pisos) & Viga & Armadura & $A_s$ (cm$^2$) & Verificación \\
\hline
20--12 & B7  & 2\,\(\phi\)22 & 7,60  & Cumple \\
20--12 & B8  & 2\,\(\phi\)22 & 7,60  & Cumple \\
20--12 & B39 & 3\,\(\phi\)22 & 11,40 & Cumple \\
\hline
11--5  & B7  & 3\,\(\phi\)22 & 11,40 & Cumple \\
11--5  & B8  & 2\,\(\phi\)22 & 7,60  & Cumple \\
11--5  & B39 & 3\,\(\phi\)22 & 11,40 & Cumple \\
\hline
4      & B7  & 2\,\(\phi\)25 & 9,82  & Cumple \\
4      & B8  & 2\,\(\phi\)25 & 9,82  & Cumple \\
4      & B39 & 3\,\(\phi\)25 & 14,73 & Cumple \\
\hline
3--1   & B7  & 2\,\(\phi\)25 & 9,82  & Cumple \\
3--1   & B8  & 2\,\(\phi\)25 & 9,82  & Cumple \\
3--1   & B39 & 3\,\(\phi\)25 & 14,73 & Cumple \\
\hline
\end{tabular}
\end{table}
 
\subsection{Diseño a corte}
 
La resistencia del hormigón al corte se calcula conforme a ACI 318-19 (sección 22.5.5.1):
\begin{equation}
    V_c = 0{,}17\,\sqrt{f_c'}\,b\,d
\end{equation}
 
Se verifica si se requiere armadura transversal:
\begin{equation}
    V_u > \tfrac{1}{2}\,\varphi\,V_c \quad \Rightarrow \quad \text{se requiere armadura de corte}
\end{equation}
 
El aporte requerido del acero transversal es:
\begin{equation}
    V_s = \frac{V_u}{\varphi} - V_c
\end{equation}
 
El espaciamiento máximo de estribos se limita según ACI 318-19:
\begin{equation}
    s_{\text{máx}} = \min\!\left(\frac{d}{2},\;600\;\text{mm}\right)
\end{equation}
y se reduce a $\min(d/4,\;300\;\text{mm})$ cuando $V_s > 0{,}33\sqrt{f_c'}\,b\,d$.
 
El área de estribo requerida y el área mínima son:
\begin{equation}
    A_{v,\text{req}} = \frac{V_s\,s}{f_y\,d}
    \qquad
    A_{v,\text{mín}} = \max\!\left(\frac{0{,}062\,\sqrt{f_c'}}{f_y}\,b\,s\;,\;\frac{0{,}35}{f_y}\,b\,s\right)
\end{equation}
 
La resistencia de diseño a corte se verifica como:
\begin{equation}
    \varphi V_n = \varphi\,(V_c + V_s) \geq V_u \qquad \checkmark
\end{equation}
 
Los resultados se presentan en la Tabla~\ref{tab:viga_corte}. En todos los casos la resistencia de
diseño supera la demanda y el área de estribo provista es mayor que la mínima requerida.
 
\begin{table}[H]
\centering
\caption{Armadura transversal adoptada por zona del edificio.}
\label{tab:viga_corte}
\begin{tabular}{llccc}
\hline
Zona (pisos) & Viga & Estribos & $s$ (m) & Verificación \\
\hline
20--1  & B7  & \(\phi\)8,  1 rama  & 0,15 & Cumple \\
20--1  & B8  & \(\phi\)8,  1 rama  & 0,15 & Cumple \\
20--12 & B39 & \(\phi\)8,  1 rama  & 0,30 & Cumple \\
11--5  & B39 & \(\phi\)8,  1 rama  & 0,30 & Cumple \\
4--3   & B39 & \(\phi\)10, 1 rama  & 0,30 & Cumple \\
2--1   & B39 & \(\phi\)12, 1 rama  & 0,30 & Cumple \\
\hline
\end{tabular}
\end{table}
 
\subsection{Resumen del diseño}
 
La Tabla~\ref{tab:viga_resumen} consolida la armadura definitiva para las tres vigas sísmicas,
diferenciando entre la zona de pisos altos y la zona de pisos bajos. Todos los elementos verifican
resistencia a flexión y a corte conforme a ACI 318-19.
 
\begin{table}[H]
\centering
\caption{Resumen de enfierradura de vigas sísmicas.}
\label{tab:viga_resumen}
\begin{tabular}{llcc}
\hline
Viga & Zona (pisos) & Armadura longitudinal & Estribos \\
\hline
B7  & 20--12 & 2\,\(\phi\)22 & \(\phi\)8 @ 0,15 m \\
B7  & 11--5  & 3\,\(\phi\)22 & \(\phi\)8 @ 0,15 m \\
B7  & 4--1   & 2\,\(\phi\)25 & \(\phi\)8 @ 0,15 m \\
\hline
B8  & 20--5  & 2\,\(\phi\)22 & \(\phi\)8 @ 0,15 m \\
B8  & 4--1   & 2\,\(\phi\)25 & \(\phi\)8 @ 0,15 m \\
\hline
B39 & 20--12 & 3\,\(\phi\)22 & \(\phi\)8  @ 0,30 m \\
B39 & 11--5  & 3\,\(\phi\)22 & \(\phi\)8  @ 0,30 m \\
B39 & 4--3   & 3\,\(\phi\)25 & \(\phi\)10 @ 0,30 m \\
B39 & 2--1   & 3\,\(\phi\)25 & \(\phi\)12 @ 0,30 m \\
\hline
\end{tabular}
\end{table}
 